\section{Discussion}
\label{sec:discussion}

The additive equation can hide persistent kinetics from bulk count,
mass, and mean size while separating number and mass sampling almost
completely. Fractional moments quantify this distinction and also
control a different object: the accumulated nonlinear displacement
of an auxiliary number path. Its finite log correction explains why
fragmentation governs the large-scale logarithmic limits and the
late-time Fourier quotient. Its almost-sure finite coagulation count
is compatible with infinite expected activity because increasingly
rare paths carry large event counts. It is not a statement that the
physical population eventually stops coagulating.

Several qualifications are part of the mathematical conclusions.
The moment coefficients are sharp only as instantaneous uniform
constants; the Golovin--Borel solution and the archived critical
simulations decay several times faster, so the estimates are Lyapunov
bounds rather than rates. The equal-split daughter envelopes and
overlap constants are extremal over their stated classes, but the
critical calendar-time exponent of the last-coagulation tail is not
determined. The envelopes reduce that exponent to the decay rate of
the deterministic number--mass overlap, a question about the equation
alone, and the simulations suggest the value $b/2$. Proving that rate,
or any asymptotic rate for the fractional moments at criticality, is
the main open problem this paper leaves.
Parent-dependent daughters retain the controlled path bounds, but
do not yield the independent translation-invariant fragmentation
reference without additional structure.

The finite-population statement shows that checking count alone can
miss a nearly maximal error in the material distribution. Its particle
input is only the material ceiling of a finite vessel, so it is a
corollary of the moment bound. It gives a sufficient logarithmic
horizon for this failure, and does not identify the first failure time or prove accuracy
before it. The exact count
law, its later diffusion limit, and the retained simulations answer
different questions from propagation of a full size distribution.
A growing-time approximation for specific bounded statistics would
require a separate argument adapted to those statistics.

The inverse application likewise separates exact identifiability
from stable estimation. A nonvanishing local amplitude cancels the
unknown initial law, and one-sided support determines the jump measure
from exact frequencies on a short low-frequency interval; that step is
analytic continuation and is severely ill-posed. The empirical sampling
bound makes attenuation of the denominator explicit and provides a
sufficient balance between coagulation bias and sampling noise, at a
certified rate whose asymptotic optimality is not asserted. It neither adapts
the observation time to unknown kinetics nor supplies stable analytic
continuation. Establishing recovery of an entire daughter law from noisy
data requires additional assumptions and regularization. Connecting
such an estimator to measurements in a finite vessel also requires an
observation model that accounts for particle dependence and continuum
error. These are distinct tasks beyond the factorization proved here.

The preparation algebra in the supplementary note offers a
complementary lesson:
what can be identified depends on which preparations and observables
are available. Preserving material loading does not remove the complete
count-neutral family. Same-shape dilution can separate linear and
quadratic initial-rate terms when coefficients remain unchanged, but
an unknown source leaves an additional ambiguity with only two
concentrations. Distributional information, physical preparation
constraints, and the sampling law must therefore be specified together.
